\subsection{\problemblock{*Fly} Data Block}\label{sec:block-fly}

TAOS uses input guidance rules to determine values for four control variables so a
trajectory can be calculated. Three of the controls specify body attitude---for
example, angle of attack, sideslip, and bank angle---and the fourth controls thrust
level through the power setting. Guidance rules and control-variable values are
entered with \problemblock{*fly} data blocks.

The \problemblock{*fly} blocks are optional. If none are given for a segment, all
control variables are set to zero. Unless unusual aerodynamic tables are supplied,
this represents a ballistic trajectory. Each \problemblock{*fly} block specifies one
guidance rule that determines one or more control-variable values. Because there are
four control variables, no more than four \problemblock{*fly} blocks can be used in
one segment. A control variable whose desired value is zero need not be specified,
because zero is the default.

\taossubhead{Body-Attitude Guidance}

The three body-attitude controls are treated separately from power setting. They can
be specified directly relative to the surrounding air mass with aerodynamic angles;
directly with conventional Euler yaw, pitch, and roll angles relative to inertial,
geocentric, or geodetic coordinates; or indirectly by supplying a desired flight
condition, such as level flight, and allowing TAOS to solve for the control values
that produce it.

The body-attitude angles that may be specified directly are listed in
\cref{tab:body-attitude-guidance-angles}. These angles combine into the nine
consistent sets shown in \cref{tab:consistent-body-attitude-sets}. All directly
specified angles in a segment must be members of one consistent set.

\begin{longtable}{@{}>{\ttfamily}l>{$}l<{$}p{0.50\linewidth}@{}}
  \caption{Body-attitude angles for vehicle guidance.}\label{tab:body-attitude-guidance-angles}\\
  \toprule
  \normalfont Variable & \normalfont Symbol & \normalfont Description \\
  \midrule
  \endfirsthead
  \toprule
  \normalfont Variable & \normalfont Symbol & \normalfont Description \\
  \midrule
  \endhead
  alpha   & \alpha        & Angle of attack \\
  alphat  & \alpha_T      & Total angle of attack \\
  bankgc  & \mu_{gc}      & Geocentric bank angle \\
  bankgd  & \mu_{gd}      & Geodetic bank angle \\
  beta    & \beta         & Angle of sideslip \\
  betae   & \beta_E       & Euler angle of sideslip \\
  phi     & \phi_w        & Windward meridian \\
  pitchgc & \Theta_{gc}   & Geocentric pitch angle \\
  pitchgd & \Theta_{gd}   & Geodetic pitch angle \\
  pitchi  & \Theta_i      & Inertial-platform pitch angle \\
  rollgc  & \Phi_{gc}     & Geocentric roll angle \\
  rollgd  & \Phi_{gd}     & Geodetic roll angle \\
  rolli   & \Phi_i        & Inertial-platform roll angle \\
  yawgc   & \Psi_{gc}     & Geocentric yaw angle \\
  yawgd   & \Psi_{gd}     & Geodetic yaw angle \\
  yawi    & \Psi_i        & Inertial-platform yaw angle \\
  \bottomrule
\end{longtable}

\begin{table}[htbp]
  \centering
  \caption{Consistent sets of body-attitude angles.}
  \label{tab:consistent-body-attitude-sets}
  \begin{tabular}{@{}clll@{}}
    \toprule
    Set & \multicolumn{3}{c}{Body-attitude angles} \\
    \midrule
    1 & \texttt{alpha}  & \texttt{betae}   & \texttt{bankgc} \\
    2 & \texttt{alpha}  & \texttt{beta}    & \texttt{bankgc} \\
    3 & \texttt{alphat} & \texttt{phi}     & \texttt{bankgc} \\
    4 & \texttt{alpha}  & \texttt{betae}   & \texttt{bankgd} \\
    5 & \texttt{alpha}  & \texttt{beta}    & \texttt{bankgd} \\
    6 & \texttt{alphat} & \texttt{phi}     & \texttt{bankgd} \\
    7 & \texttt{yawgc}  & \texttt{pitchgc} & \texttt{rollgc} \\
    8 & \texttt{yawgd}  & \texttt{pitchgd} & \texttt{rollgd} \\
    9 & \texttt{yawi}   & \texttt{pitchi}  & \texttt{rolli} \\
    \bottomrule
  \end{tabular}
\end{table}

Often all three angles in a set are entered directly in three separate
\problemblock{*fly} blocks, but this is not required. If only one or two angles are
specified, the remaining angles are free. Other guidance rules may determine those
free angles, or they remain at their default value of zero.

The usual form of a \problemblock{*fly} block is the keyword followed by a
guidance variable and its value. For example,

\begin{lstlisting}[style=taosmanual]
*fly alpha=5.0
*fly betae=2.5
\end{lstlisting}

requests an angle of attack of $5^\circ$ and an Euler sideslip of $2.5^\circ$. The
specified variables appear in both Set 1 and Set 4, but TAOS chooses Set 1 because it
has the lower set number. The remaining body-attitude variable is therefore
\texttt{bankgc}. Because it is not specified and no other guidance rule determines
it, geocentric bank angle defaults to zero. Power setting also defaults to zero.

Use of \texttt{alpha} and \texttt{betae} is encouraged over \texttt{alpha} and
\texttt{beta}. The projection-based sideslip angle \texttt{beta} becomes awkward
when angle of attack approaches $\pm90^\circ$; TAOS modifies the standard
definition to avoid a singularity, but the \texttt{alpha}/\texttt{betae} combination
is the more robust representation (\cref{sec:wind-coordinate}).

The following pair is inconsistent and generates an input error because total angle
of attack is paired with windward meridian, not Euler sideslip:

\begin{lstlisting}[style=taosmanual]
*fly alphat=5.0       # Incorrect: inconsistent angle set
*fly betae=5.0
\end{lstlisting}

\taossubhead{Referencing Current Values}

An asterisk can be entered in place of a numerical value to request the value at the
end of the preceding segment:

\begin{lstlisting}[style=taosmanual]
*fly yawgd=*
*fly pitchgd=*
*fly rollgd=*
\end{lstlisting}

Asterisk values should not be used in the first segment of a trajectory because no
preceding values exist.

\taossubhead{Power Setting}

The \statevar{power} variable is a special control that can be used with propulsion
tables to simulate an engine throttle. If the thrust and mass-flow tables are functions
of \statevar{power}, changing it changes thrust and fuel flow. Power setting is
nondimensional; its numerical scale can represent engine speed, throttle position, or
another convenient quantity, provided the same convention is used in the table and
problem files. It may be specified directly or determined indirectly from a desired
flight condition such as constant Mach number.

\taossubhead{Flight-Condition Guidance}

Control variables can also be specified indirectly with desired flight conditions. In
this case TAOS identifies the controls that have not been directly specified---the
free control variables---and varies them with the iterative Newton--Raphson method
described in \cref{sec:control-variable-solution} until the requested conditions are
satisfied. The available flight-condition rules are listed in
\cref{tab:flight-condition-guidance}.

\begin{longtable}{@{}>{\ttfamily}lp{0.69\linewidth}@{}}
  \caption{Flight conditions for vehicle guidance.}\label{tab:flight-condition-guidance}\\
  \toprule
  \normalfont Variable & Description \\
  \midrule
  \endfirsthead
  \toprule
  \normalfont Variable & Description \\
  \midrule
  \endhead
  alt       & Fly at a constant altitude \\
  cl        & Fly at a constant lift coefficient \\
  cs        & Fly at a constant side-force coefficient \\
  downria   & Fly down the range-insensitive axis \\
  dynprs    & Fly at a constant dynamic pressure \\
  gamgc     & Fly at a constant geocentric flight-path angle \\
  gamgd     & Fly at a constant geodetic flight-path angle \\
  intercept & Fly to intercept another trajectory \\
  l/d       & Fly at a constant lift-to-drag ratio \\
  l/d-max   & Fly at the maximum lift-to-drag ratio \\
  mach      & Fly at a constant Mach number \\
  nx        & Fly at a constant axial specific-load factor \\
  ny        & Fly at a constant lateral specific-load factor \\
  nz        & Fly at a constant normal specific-load factor \\
  propnav   & Fly to intercept with proportional navigation \\
  psigc     & Fly at a constant geocentric heading angle \\
  psigd     & Fly at a constant geodetic heading angle \\
  thrust    & Fly at a constant thrust \\
  upria     & Fly up the range-insensitive axis \\
  vel       & Fly at a constant velocity \\
  \bottomrule
\end{longtable}

The search normally converges, but convergence is not guaranteed. If a problem
occurs, first verify that the free controls can actually influence the requested flight
conditions. For example, with zero bank angle, angle of attack has no effect on
heading and cannot be used to solve for a commanded heading. Next inspect the
control limits in the \problemblock{*limits} block. Restricting free angles to the
range covered by the aerodynamic tables prevents unreasonable extrapolation. Finally,
inspect \statevar{dtguid} in the \problemblock{*integ} block. This value acts as a
time constant or gain controlling how rapidly the trajectory is corrected from the
actual state to the desired state; a value that is either too large or too small can
prevent the trajectory from approaching the command.

The following set of rules requests level flight at Mach 0.8:

\begin{lstlisting}[style=taosmanual]
*fly gamgd=0
*fly mach=0.80
*fly bankgd=0
*fly betae=0
\end{lstlisting}

Two controls are fixed directly, \texttt{bankgd} and \texttt{betae}, and two are
specified indirectly. TAOS solves for angle of attack and power setting to obtain a
zero geodetic flight-path angle and Mach 0.8. The last two blocks are optional because
both variables default to zero; they are shown for clarity.

\taossubhead{Special Guidance Rules}

The rules \texttt{l/d-max}, \texttt{intercept}, \texttt{propnav},
\texttt{downria}, and \texttt{upria} require special handling. The maximum
lift-to-drag rule solves for the angle of attack that maximizes $L/D$ and requires no
value:

\begin{lstlisting}[style=taosmanual]
*fly l/d-max
\end{lstlisting}

The predictive-intercept rule takes the trajectory number of the target:

\begin{lstlisting}[style=taosmanual]
*fly intercept=3
\end{lstlisting}

This directs the vehicle toward the estimated intercept point using the method in
\cref{sec:intercepts-pn}. Proportional navigation uses the same
form:

\begin{lstlisting}[style=taosmanual]
*fly propnav=2
\end{lstlisting}

and uses \statevar{dtguid} from the \problemblock{*integ} block as the navigation
constant. Both intercept rules require one free control for yaw and one for pitch, so
each is equivalent to two ordinary guidance rules. When either is used, a segment
can contain at most three \problemblock{*fly} blocks.

The \texttt{downria} and \texttt{upria} rules align the body $x$ axis with the range-
insensitive axis. A velocity increment along this axis changes time to impact without
changing impact range. The former selects the direction that reduces flight time and
the latter the direction that increases it (\cref{sec:range-insensitive-axis}). Their
value is the impact altitude used by the initial-impact-point calculation, usually zero:

\begin{lstlisting}[style=taosmanual]
*fly downria=100000
\end{lstlisting}

This example solves for the range-insensitive axis using an assumed impact altitude
of 100,000 ft. Like the intercept rules, the range-insensitive rules require free yaw
and pitch controls and are equivalent to two ordinary guidance rules.

\taossubhead{Rate Guidance}

Body-attitude rates and flight-condition rates are specified by appending
\texttt{dt}---for ``dot,'' or first derivative with respect to time---to a guidance
variable. For example, the following rules request level flight, a geodetic heading
rate of $5~\mathrm{deg/s}$, zero Euler sideslip, and a power-setting rate of 1000
units/s:

\begin{lstlisting}[style=taosmanual]
*fly gamgd=0
*fly psigddt=5.0
*fly betae=0
*fly powerdt=1000
\end{lstlisting}

If power setting is proportional to thrust in the propulsion tables, the result is a
level, accelerating turn. TAOS varies angle of attack and geocentric bank angle to
produce the requested heading rate while maintaining level flight.

\taossubhead{Guidance Tables}

If constants and rates do not describe the desired trajectory, a guidance value can be
interpolated from a table as a function of a state variable:

\begin{lstlisting}[style=taosmanual]
*fly alpha vrs tseg interp-2
  0.0   0.0
  5.0   1.0
  6.0   2.0
  5.0   3.0
  0.0   4.0
\end{lstlisting}

The required keyword \texttt{vrs} separates the guidance variable from the
independent state variable. Guidance tables may depend on the same state variables
as table-file data (\cref{tab:taos-state-variables}). Three interpolation methods
are available:

\begin{center}
\begin{tabular}{@{}>{\ttfamily}ll@{}}
  interp-1 & Linear interpolation (default) \\
  interp-2 & Circular interpolation \\
  interp-3 & Polynomial interpolation \\
\end{tabular}
\end{center}

Guidance tables are used especially in flight-path optimization, where an angle-of-
attack or other control history is varied to optimize a flight condition subject to
constraints. During optimization, table values are perturbed by very small amounts
when numerical gradients are formed. To evaluate those gradients accurately, TAOS
must integrate exactly to each point in a time-based guidance table. It therefore
adjusts the integration step automatically so that every point in a table depending on
\statevar{time}, \statevar{tseg}, or \statevar{tmark} lies at the end of an integration
step.
